\chapter{\IfLanguageName{dutch}{Stand van zaken}{State of the art}}%
\label{ch:stand-van-zaken}

% Eigen stand van zaken: 
\section{OWASP en hedendaagse web- en API-dreigingen}
\label{sec:api-security-context}
Om de rol van anomaliedetectie binnen webbeveiliging te kaderen, is het belangrijk eerst de huidige dreigingen te schetsen en de beperkingen van bestaande detectiemechanismen te bespreken. 

Binnen web- en API-beveiliging speelt OWASP, het Open Web Application Security Project, een centrale rol voor het identificeren van de tien meest kritieke kwetsbaarheden. OWASP publiceert periodiek top tien lijsten voor zowel webapplicaties als voor API's\autocite{OWASP2026,OWASP2023}. Deze lijsten bieden een overzicht van de meest voorkomende en impactvolle kwetsbaarheden binnen het domein van webbeveiliging.%OWASP WEB EN API BRON. 

Voor webbapplicaties omvat de OWASP TOP 10 onder meer kwetsbaarheden zoals injection-aanvallen, broken authentication en security misconfigurations. Deze kwetsbaarheden zijn voornamelijk gericht tot de interactie tussen gebruiker en de webinterface\autocite{OWASP2026}%OWASP WEB BRON. 

Voor API's verschuift de focus naar andere soorten risico's zoals Broken Object Level Authorization (BOLA), Broken Function Level Authorization en Unrestricted Resource Consumption. API kwetsbaarheden zijn vaker gerelateerd aan onvoldoende controle op de toegang tot deze APIs en het misbruiken van businesslogica\autocite{OWASP2023}%API BRON. 

Het onderscheid dat gemaakt wordt tussen web- en API-kwetsbaarheden benadrukt dat er nood is aan beveiliging op meerdere niveaus. Waar dat webaanvallen eerder gericht zijn op de gebruikersinterface, richten API-aanvallen zich tot de directe communicatie tussen systemen. Dit onderscheid is de basis voor een diepere analyse van hoe aanvallen op beide niveaus evolueren. 

\subsection{De groeiende complexiteit van Web-aanvallen}
\label{subsec:api-aanvallen-complexiteit}
De complexiteit van aanvallen gericht op webapplicaties is in de afgelopen jaren sterk toegenomen. Volgens rapporten van Akamai is er sprake een significante stijging in webaanvallen, hier spreken ze onder andere over een toename van 33 procent tussen 2023 en 2024\autocite{Akamai2025}%Akamai bron.

Binnen de OWASP TOP 10 komen verschillende aanvalsvormen naar voren die deze evolutie illustreren. Een belangrijk voorbeeld is Broken Access Control waar een aanvaller buiten zijn toegestane rechten handelingen uitvoert. Daarnaast blijven injection attacks een voorkomend risico, dit is wanneer de aanvaller code zoals SQL of OS commando's via de webapplicatie uit laat voeren op de server waar de applicatie gehost word om zo data te verkrijgen. Ook authentication failures vormen een ketsbaarheid, dit is wanneer de beveiliging van de applicatie er onvoldoende in slaagt om de user van de aanvaller te detecteren als ongeldig of incorrect dit kan dan leiden tot brute force en geautomatiseerde aanvallen\autocite{OWASP2026}%bron owasp WEB.

Een kritieke evolutie binnen deze aanvalsvormen is dat kwaadaardig verkeer steeds moeilijker te onderscheiden is van normaal geldig verkeer. Vandaag de dag maken aanvallers steeds vaker gebruik van geautomatiseerde bots. Deze bots zijn speciaal ontworpen om menselijk gedrag zo nauwkeurig mogelijk te simuleren door het implementeren van variatie in de frequentie van verzonden requests en hun navigatiepatronen. Deze geautomatiseerde aanvallen maken gebruik van grote aantallen requests die geldig lijken en zich aan kunnen passen naar de beveiliging van een webapplicatie. Dit zorgt ervoor dat het verschil tussen normaal en kwaadaardig verkeer moeilijker waarneembaar wordt\autocite{F52026}.%F5 bron 

Naast de stijgende complexiteit van webaanvallen kent ook het API-landschap een gelijkaardige evolutie. 
\subsection{De groeiende complexiteit van API-aanvallen}
\label{subsec:api-aanvallen-complexiteit}

De complexiteit en frequentie van aanvallen op APIs is eveneens sterk toegenomen. APIs vormen een essentieel onderdeel voor de moderne webapplicaties, maar ze vergroten het aanvalsoppervlak aanzienlijk\autocite{Akamai2025}. %Akamai bron

In tegenstelling tot klassieke webaanvallen richten API-aanvallen zich op de communicatie tussen softwarecomponenten. Volgens OWASP ontstaan API-kwetsbaarheden voornamelijk door foutieve of beperkte authenticatie en autorisatie. Kwetsbaarheden zoals Broken Object Level Authorization en Broken Function Level Authorization tonen aan dat aanvallers zich richten tot endpoints die technisch correct functioneren, maar onvoldoende toegangscontrole bevatten.

Daarnaast worden API’s vaak misbruikt via geautomatiseerde requests, bijvoorbeeld bij Unrestricted Resource Consumption en business logic abuse. Dit type aanvallen maakt gebruik van legitieme interacties, waardoor detectie via traditionele signature-based systemen moeilijk is\autocite{OWASP2023}.%bron OWASP 

Zowel bij webaanvallen als bij API-aanvallen stelt zich dezelfde vraag: waarom slagen bestaande beveiligingsmechanismen er onvoldoende in om deze dreigingen tijdig te detecteren. 

\subsection{Beperkingen van traditionele detectiemechanismen}
\label{subsec:traditioneel-beperkingen}
% Leg uit waarom rule-based systemen tekorschieten bij moderne dreigingen.

Traditionele beveiligingsmechanismen zoals Web Application Firewall (WAF), signature based Intrusion Detection Systems (IDS) en andere regelgebaseerde detectiesystemen steunen voornamelijk op bekende aanvalssignaturen en patronen die op voorhand gedefineerd zijn. Zoals \textcite{Standards2007} beschrijft in hun Intrusion Detection and Prevention systems richtlijnen zijn deze systemen zeer effectief in het detecteren van gedocumenteerde dreigingen, maar eerder beperkt in het detecteren van nieuwe aanvallen waarvoor nog geen signature aanval bestaat. 

Zoals eerder beschreven maken moderne aanvallers steeds vaker gebruik van geautomatiseerde technieken die gespecialiseerd zijn in het nabootsen van legitiem verkeer en dat zich aanpast aan de beveiliging van een applicatie. Dergelijke aanvallen bevatten vaak geen gekende signatuur of vertonen geen afwijkend patroon dat op voorhand gedifinieerd kan worden \autocite{F52026}.

Recente literatuur bevestigt deze claims en benadrukt signature-based systemen te beperkt zijn in het detecteren van zero-day aanvallen binnen het snel evoluerende cybersecuritylandschap. Daarom wordt machine learning steeds vaker voorgesteld als een aanvulling op deze bestaande systemen om eerder ongekende aanvallen te detecteren, false positives te verminderen en beveiliging te verbeteren\autocite{Kumar_2025}. % komt van deze bron Impact of Machine Learning on Intrusion Detection Systems for the Protection of Critical Infrastructure


\section{Anomaliedetectie binnen Security Operations}
\label{sec:anomaliedetectie-soc}
%Zou hier graag nog een subsectie aan toe willen voegen

De beperkingen van traditionele detectiemechanismen, zoals beschreven in \ref{subsec:traditioneel-beperkingen}, creëren een behoefte voor een aanvullend detectiesysteem binnen het Security Operations Team.  Binnen de context van web- en API-beveiliging speelt het SOC een centrale rol in het bewaken van inkomend verkeer en het detecteren van verdachte events.

\subsection{Akamai APP en API-protector}
\label{subsec:App-API}
Akamai App \& API Protector is een geïntegreerde beveiligingsoplossing die webbapplicaties en API's beschermt tegen een breed aantal dreigingen, waaronder Dos-aanvallen, botactiviteit en gekende kwetsbaarheden uit onder andere de OWASP top tien. Dit doet het door het combineren van meerdere modules waaronder een WAF, botmitigatie en API-bescherming in één platform.%https://www.akamai.com/products/app-and-api-protector

De bescherming van webapplicaties gebeurt via een regelgebaseerde evaluatie van inkomende requests. Indien een request voldoet aan vooraf gedefinieerde criteria, past het platform een set beveiligingscontroles toe. Indien een aanvalsvector wordt gedetecteerd, wordt de toegang geblokkeerd of wordt er een alert gegenereerd afhankelijk van de geconfigureerde instellingen\autocite{Technologies2026}.

Voor botdetectie maakt het platform gebruik van Bot-manager-module, die naast regelgebaseerde signaturen ook machine learning inzet voor het detecteren van bots. Deze module richt zich op een breed spectrum aan detectiemethoden: browser fingerprinting, request-anomalieën, reputatiescoring en gedragsanalyse. De Bot manager is specifiek ontworpen voor het indentificeren van bots binnen het verkeer tussen de menselijke gebruikers.%https://www.akamai.com/products/bot-manager 

Voor API-bescherming biedt het platform zichtbaarheid over alle API-endpoints binnen een omgeving, inclusief niet-gedocumenteerde of zogenaamde shadow  API's. Het platform analyseert API-verkeer op basis van vooraf bepaalde beleidsregels en detecteert afwijkingen zoals buitensporige requests of ongeautoriseerde toegangspogingen\autocite{Akamai2025}.
\subsection{De toegevoegde waarde van ML in SOC-context}
\label{subsec:toegevoegde-waarde-ML}
% Motiveer waarom ML een waardevolle aanvulling is op klassieke detectiesysteem bron: https://dl.acm.org/doi/10.1145/3723158?
Een Security Operations Center (SOC) is verantwoordelijk voor het continu monitoren en analyseren van web- en API-verkeer op zoek naar verdachte activiteit en beveiligingsincidenten. Binnen een webbeveiligingscontext maken SOC-analisten gebruik van tools zoals WAF-systemen, IDS-systemen en SIEM-platformen (Security Information and Event Management) om inkomende requests te evalueren en potentiële aanvallen tijdig te detecteren en te beantwoorden.

In moderne Security Operations Teams leidt een grote hoeveelheid security alerts tot een fenomeen wat door de literatuur beschreven wordt als alert fatigue. Dit is wanneer analisten worden overweldigd door een constante stroom van alerts en false positives. Met als gevolg een verhoogde operationele belasting, wat het risico biedt dat kritieke incidenten worden gemist. 

De toegevoegde waarde van machine learning binnen een SOC-team heeft voordelen op verschillende niveaus: na het inwerken van het model een reductie van false positives, verbeterde detectie van abnormaal gedrag en het schakelen van reactief handelen naar proactief detecteren\textcite{Miranda2025}. %https://www.crowdstrike.com/en-us/cybersecurity-101/next-gen-siem/ai-siem/


\section{Machine learning-methoden voor anomaliedetectie}
\label{sec:ML-methoden-anomalie}

Zoals beschreven in sectie \ref{subsec:traditioneel-beperkingen} hebben de beperkingen van traditionele rule-based detectiesystemen ervoor gezorgd dat de interesse in machine learning gedreven benaderingen voor anomaliedetectie zijn gestegen. Deze sectie biedt een overzicht van de belangrijkste machine learning gedreven benaderingen voor anomaliedetectie, met specifieke aandacht op hun toepasbaarheid. 

\subsection{supervised, unsupervised en semi-supervised}
\label{subsec:toegevoegde-waarde-ML}
% overzicht geven van de 3 hoofdbenaderingen met voor- en nadelen. 

Machine learning benaderingen voor anomaliedetectie worden onderverdeeld in 3 categorieën op basis van de nood aan gelabelde trainingsdata. \textcite{Saabith2023} maakt een onderscheid in drie hoofdcategorieën: supervised learning, unsupervised learning en semi-supervised learning. Deze benaderingen bevatten elk sterke punten en zwakke punten die hun toepasbaarheid bepalen binnen security context.


\subsubsection{Supervised learning}

Supervised learning-algoritmen vereisen een dataset waarbij elk datapunt uitdrukkelijk gelabeld is als normaal of abnormaal gedrag, vervolgens leert het model tijdens de trainingsfase de relatie tussen de invoerkenmerken (features in de dataset) en de daarbij horende labels hierdoor is het instaat om nieuwe datapunten te classificeren~\autocite{Saabith2023}. 
De algoritmen binnen deze categorie zijn: bayesian networks, k-nearest neighbors, random forest, supervised neural networks en suport vector machines. Over het algemeen ligt de detectierate hoger bij supervised modellen dan bij unsupervised modellen, maar dit gaat dan ook gepaard met de moeilijkheidsgraad in het trainen van het model waarbij dat bij supervised hoger ligt\autocite{Kumari2024}. 
De voornaamste beperking van supervised learning-algoritmenn is de afhankelijkheid aan een gelabelde dataset. In werkelijkheid zijn gelabelde datasets weinig te vinden. Ten eerste vereist het labelen van data security-analisten met expertise in het domein dat duizenden events dienen te categoriseren, een proces dat lang duurt~\autocite{Kumar_2025}. 
Vervolgens kan er niet worden gegarandeerd dat de trainingsdata alle relevante aanvalstypes bevat, zo is het mogelijk dat niet bekende aanvalstypes, zero-day exploits en andere varianten van aanvallen mogelijks niet aanwezig zijn in de trainingsdata van het model. Dit kan ervoor zorgen dat het model nieuwe subtiele anomalieën mist~\autocite{Kumari2024}.

\subsubsection{Unsupervised learning}

Unsupervised learning-algoritmen zijn in tegenstelling tot supervised benaderingen niet gebonden aan een gelabelde dataset waarop het model getraind moet worden. Deze modellen leren om patronen en structuren te herkennen in de data op basis van zowel gelijkheden als verschillen tussen datapunten. Punten die van het opgebouwde patroon afwijken of niet binnen een bestaande vooraf gedefinieerde  cluster past zullen worden gezien als een anomalie\textcite{Saabith2023}. 
Algoritmen binnen deze categorie zijn: clustering methoden (K-means, DBSCAN), isolation forest, local outlier factor, Principal Component Analysis en Autoencoders. een van de voordelen van unsupervised machine learning modellen is dat ze uitblinken in het detecteren van ongekende anomalieën doordat ze zich niet baseren op historische trainingsdata\autocite{Kumari2024}.

Een andere nuttige eigenschap van unsupervised learning is dat deze modellen continue leren wat normaal is op basis van de binnenkomende data, dit zorgt ervoor dat ze schaalbaar zijn in omgevingen waar veel data binnenkomt en ze zich kunnen aanpassen naarmate de data veranderd\autocite{Traceable2024}.

De voornaamste beperking van unsupervised learning is dat het kan resulteren in een hogere false positive rate ten opzichte van een supervised benadering. Doordat het model geen expliciete instructies ontvangen heeft zal het alle statistische datapunten die uitschieten beschouwen als verdacht, dit kan dus ook zeldzaam legitiem verkeer bevatten\autocite{Saabith2023}. 


\subsubsection{Semi supervised learning}

Semi-supervised learning is een combinatie benadering van beide voorgaande methoden, het model wordt dus getraind op zowel niet gelabelde data als gelabelde data. De unlabeld data dienen om de accuraatheid van het model te verhogen op niet gekende anomaliteiten en de labeld data geeft het model zicht op welke afwijkingen er voornamelijk zijn\autocite{Saabith2023}. 
Algoritmen binnen deze categorie zijn one-class SVM en autoencoder met training data. Autoencoder kan zowel usnupervised als semi-supervised dit is afhankelijk van het wel of niet implementeren van labeld data\autocite{Bajallan2020}. %komt uit deze bron https://kth.diva-portal.org/smash/get/diva2:1465807/FULLTEXT01.pdf
Doordat semi-supervised beperkte labeld data gebruikt kan het accuratere beslissingen maken op wat wel of niet een afwijking is op het normale gedrag. Hierdoor kan het ten opzichte van een unsupervised benadering de false positive rate reduceren\autocite{EmergentMind_2026}.  % komt uit https://www.emergentmind.com/topics/semi-supervised-anomaly-detection bron 
De voornaamste restrictie van een semi-supervised benadering is dat er hoewel in mindere maten nog steeds nood is aan gelabelde data voor het trainen van het model \autocite{Saabith2023}.

\subsection{Veelgebruikte algoritmen per categorie}%
\label{subsec:ml-algoritmen-overzicht}

\subsubsection{Supervised ML modellen}

\paragraph{Random Forest}

Random Forest is een ensemble-gebaseerde supervised machine learning algoritme dat gebruikt kan worden voor anomaliedetectie. 
Het model bestaat uit een verzameling van decision trees, waarbij elke boom getrained wordt op een willekeurige subset uit de trainingsdata en een willekeurige set van features gebruikt\autocite{Breiman2001}.

Random Forest kan herleid worden tot een verzameling boomgebaseerde classifiers, hierbij is elke boom opgebouwd door een onafhankelijke en identieke vector die willekeurig gekozen is door het model. 
Voor een gegeven inputvector zal het model een voorspelling geven, waarna de classificatie bepaald wordt door majority voting.
In de context van anomaliedetectie zal de boom een classificatie geven zoals normaal of anomaal, en de uiteindelijke beslissing is gebaseerd op een meerderheid in stemmen\autocite{Primartha2017}.

Binnen Random Forest zijn er enkele essentiële parameters die direct invloed hebben op de werking en stabiliteit van het model. 
De \texttt{n\_estimators} parameter staat in voor het aantal decision trees binnen het ensemble; een hoger aantal bomen verhoogt de stabiliteit maar heeft ook rechtstreeks invloed op de rekencomplexiteit. 
De \texttt{max\_features} parameter bepaalt hoeveel features willekeurig gekozen worden per splitsing in de decision tree\autocite{Breiman2001}.

Een fundamentele beperking van Random Forest is de nood aan gelabelde trainings data
waarbij elk datapunt in de dataset gelabeld moet zijn als normaal of anomaal. Dit maakt Random Forest minder geschikt voor omgevingen waarin de aard van het verkeer niet op voorhand gekend is en de nadruk ligt op het ontdekken van voorheen onbekende anomalieën~\autocite{Breiman2001}

\textbf{Voordelen:}
\begin{itemize}
    \item Robuust tegen overfitting, doordat het model niet sterk afstemt op
    zijn trainingsdata en goed generaliseert naar nieuwe, ongeziene data.
    \item Expliciete featureselectie is doorgaans niet noodzakelijk, doordat
    het model efficiënt omgaat met een groot aantal variabelen.
\end{itemize}

\textbf{Nadelen:}
\begin{itemize}
    \item Lage interpreteerbaarheid, doordat de uiteindelijke voorspelling
    gemaakt wordt door een combinatie van decision trees, is het moeilijk om
    de grondslag van een classificatie te achterhalen.
    \item Prestaties kunnen dalen bij sterk gecorreleerde features, wat de
    diversiteit tussen de decision trees vermindert.
\end{itemize}

\medskip

\paragraph{Support Vector Machine}

Support Vector Machine is een supervised classificatiealgoritme dat een scheidingslijn genaamd de hyperplane trekt tussen twee klassen. 
Bij binaire classificatie zoals het onderscheid maken tussen normaal en anomalie, probeert SVM de grens te vinden die de twee klassen scheidt met een zo groot mogelijke marge. 
De marge is de afstand tussen de beslissingsgrens en de datapunten van beide klassen die zich het dichtst bij bevinden, deze datapunten worden ook welde support vectors genoemd. 
Door de marge zo groot mogelijk te maken kan SVM beter generaliseren naar nieuwe data en is het minder gevoelig aan kleine variaties in de set\autocite{Cortes1995}.

In tegenstelling tot de praktijk is hier data zelden linair scheidbaar: de grens tussen normaal of anomaal kan zelden met een rechte lijn worden getrokken. 
Om deze uitdaging aan te pakken maakt SVM gebruik van kernelfuncties. Een kernelfunctie transformeert data richting een hogere dimensionale ruimte waarin de data dan weer wel linair scheidbaar is, dit wordt ook wel de kernel trick genoemd.
Het gebruik van een kernelfunctie heeft ook een rechtstreeks positieve invloed op de berekening van de transformatie. 
SVM gebruikt voornamelijk kernelfuncties zoals de lineaire kernel, de polynomiale kernel en de Radial Basis Function kernel. De keuze van kernel is sterk afhankelijk met de structuur van de data. %BRON: 

\textbf{Voordelen:}
\begin{itemize}
    \item Sterk in hoog-dimensionale datasets, aangezien het model goed overweg
    kan met een groot aantal features binnen de data.
    \item Capabel model voor het aanduiden van grenzen tussen klassen in de data.
\end{itemize}

\textbf{Nadelen:}
\begin{itemize}
    \item Lage interpreteerbaarheid, het is moeilijk om te achterhalen wat de
    grondslag was van een bepaalde classificatie.
    \item Gevoelig aan parameter settings en trainingstijd kan exponentieel
    stijgen bij grote datasets.
\end{itemize}

\medskip

\subsubsection{Unsupervised ML modellen}

\paragraph{Isolation Forest}

Isolation Forest is een unsupervised anomaliedetectiemethode dat vertrekt vanuit een andere invalshoek ten opzichte van andere ML modellen: in plaats van een profiel op te bouwen van hoe normaal gedrag eruitziet, focust Isolation Forest zich op het isoleren van afwijkende datapunten. 
Dit is mogelijk doordat het model ervan uitgaat dat anomalieën doorgaans zeldzaam zijn en afwijkende waarden bezitten\autocite{Liu2008}.

Het algoritme bouwt een vooraf gedefinieerd aantal binaire decision trees. 
In elke boom wordt iteratief een willekeurige feature en een willekeurige splitswaarde geselecteerd om datapunten te scheiden. 
De nodige padlengte om een datapunt volledig te isoleren wordt gemeten en genormaliseerd tot een anomaliescore tussen 0 en -1.
Datapunten met een korte gemiddelde padlengte over alle bomen heen krijgen een hoge anomaliescore\autocite{Lesouple2021}.

Isolation Forest maakt gebruik van subsampeling voor het trainen van de decision trees binnen het model, dit wil zeggen dat niet elke boom getrained wordt op de volledige dataset maar op een willekeurige subset van de data.
Dit verbetert de detectiekwaliteit van het model doordat het 2 fenomenen tegengaat. 
Enerzijds voorkomt het masking, dit is wanneer datapunten die dicht bij elkaar liggen in een grote dataset onterecht door het model zouden worden beschouwd als een kleine groep legitieme data. 
Anderzijds voorkomt het swamping, dit is het tegenovergestelde van masking, hierbij worden datapunten die zicht dicht tegen de rand bevinden verward als anomaal. 
Een ander voordeel van subsampeling is dat het model efficiënt schaalt naar grote datasets doordat de trainingstijd per boom beperkt blijft. 
Bij de training van het model dient een contamination-parameter te worden gedefineerd. Deze parameter is verantwoordelijk voor het bepalen welk percentage van de data het model zal classificeren als afwijkend.
De keuze van deze parameter is afhankelijk van de context van het model en de data en heeft rechtstreekt invloed op de balans tussen het detecteren van normaal of anomaal verkeer\autocite{Liu2008}. 

\textbf{Voordelen:}
\begin{itemize}
    \item Zeer efficiënt en schaalbaar, doordat het algoritme een lineaire
    tijdcomplexiteit heeft en weinig geheugen vereist, kan het goed toegepast
    worden op grote datasets en data met veel dimensies.
    \item Effectief in het detecteren van anomalieën, omdat het model anomalieën
    direct probeert te isoleren in plaats van normale data te modelleren,
    waardoor afwijkende punten sneller herkend worden.
\end{itemize}

\textbf{Nadelen:}
\begin{itemize}
    \item Moeilijkheden bij overlappende of dichte clusters, wanneer anomalieën
    dicht bij normale data liggen of in grote groepen voorkomen, worden ze minder
    goed geïsoleerd en dus moeilijker gedetecteerd.
    \item Gevoelig voor data-eigenschappen zoals dimensionaliteit en sample size,
    waardoor prestaties kunnen afnemen bij veel irrelevante features of een
    ongunstige keuze van de subsample grootte.
\end{itemize}

\medskip

\paragraph{Local Outlier Factor}

Local Outlier Factor is een unsupervised anomaliedetectiemethode die afwijkende datapunten detecteert op basis van dichtheidsverschillen. LOF evalueert een datapunt relatief ten opzichte van zijn k-dichtste buren.

Het algoritme start met het berekenen van de \texttt{k-distance}, dit is de afstand tussen het datapunt en de k-de dichtste buur. 
Vervolgens wordt de reachability distance berekend, een metric die instabiliteit in de meting voorkomt zodat twee punten die dicht bij elkaar liggen niet automatisch als normaal worden beschouwd. 
De uiteindelijke LOF-score is de verhouding tussen de gemiddelde lokale dichtheid van de buren van het datapunt en de lokale dichtheid van het punt zelf. 
Wanneer de LOF-waarde zakt onder 1 wijst dit op een lagere dichtheid ten opzichte van andere punten en dus een mogelijke anomalie\autocite{Mavuduru2026}.

De voornaamste parameter van LOF is \texttt{k}, het aantal buren dat het model evalueert. 
Een kleinere k-waarde maakt het model gevoeliger voor lokale variaties, een grotere k-waarde maakt het model beter bestendig tegen false positives maar ook minder gevoelig voor lokale anomalieën.

Het voornaamste verschil met isolation forest is dat LOF anomalieën detecteert op basis van relatieve dichtheidsverschillen, dit maakt het model effectiever in het detecteren van lokale anomalieën die zich dichter bevinden bij normale clusters. 
De berekining van de k-nearest neighbors voor elk datapunt schaalt kwadratisch waardoor het minder geschikt is voor grote datasets in een real-time omgeving~\autocite{Breunig2000}. 

\textbf{Voordelen:}
\begin{itemize}
    \item Effectief in het detecteren van lokale anomalieën, doordat het model
    rekening houdt met de dichtheid van naburige datapunten en afwijkingen
    identificeert ten opzichte van de directe omgeving.
    \item Geen gelabelde data vereist, aangezien LOF een unsupervised algoritme
    is dat anomalieën detecteert op basis van structurele eigenschappen van de
    dataset.
\end{itemize}

\textbf{Nadelen:}
\begin{itemize}
    \item Gevoelig voor parameterkeuze, in het bijzonder het aantal buren
    (\texttt{k}), wat een grote invloed heeft op de prestaties en kan leiden
    tot foutieve classificaties.
    \item Kans op false positives aan de randen van clusters, doordat punten
    met een lagere dichtheid aan de grens van normale clusters ten onrechte als
    anomalieën kunnen worden beschouwd.
\end{itemize}

\medskip

\subsubsection{Semi-supervised ML modellen}

\paragraph{One-Class SVM}

One-Class SVM is een variant op de klassieke SVM waarbij het model uitsluitend
getraind wordt op data die als normaal wordt beschouwd. Het doel is om een
beslissingsgrens te leren die de normale data zo goed mogelijk omvat en
afwijkende punten uitsluit.

In tegenstelling tot de klassieke SVM, die twee klassen van elkaar scheidt, zal OC-SVM proberen om een hyperplane te vinden die de data scheidt van de oorsprong.
Doordat datapunten niet altijd direct scheidbaar zijn, maakt SVM gebruik van een kernelfunctie die data projecteert naar een hogere dimensionale ruimte waardoor scheiding mogelijk wordt. 
Het model gebruikt een parameter die de balans bepaalt tussen de toegestane outliers en de complexiteit van het model: enerzijds bepaalt deze een bovengrens voor het aantal toegestane anomalieën, anderzijds een ondergrens voor de toegestane hoeveelheid support vectors\autocite{Hejazi2013}.

De belangrijkste hyperparameter bij one-class SVM is de $\nu$-parameter (\textit{nu}), een waarde tussen 0 en 1 met 2 belangrijke functies. 
Het dient als een bovengrens op het onderdeel aan trainingsfouten, dit zijn de outliers binnen de trainingsdata, en eveneens als een ondergrens op het onderdeel van de support vectors.
Een lagere $\nu$-waarde kan resulteren in een striktere beslissingsgrens die minder outliers toestaat, terwijl een hogere waarde eerder het effect kan hebben dat uitschieters worden geaccepteerd~\autocite{Scholkopf2001}. 

\textbf{Voordelen:}
\begin{itemize}
    \item Geschikt voor anomaliedetectie met één enkele klasse, doordat het model
    een grens leert rond normaal gedrag zonder nood aan expliciet gelabelde
    anomalieën.
    \item Goede prestaties bij imbalanced datasets, aangezien het model ontworpen
    is om zeldzame afwijkingen te detecteren en een lage generalization error kan
    bereiken bij correcte parameterinstelling.
\end{itemize}

\textbf{Nadelen:}
\begin{itemize}
    \item Sterke afhankelijkheid van hyperparameters en kernelkeuze, waarbij een
    verkeerde configuratie kan leiden tot suboptimale prestaties.
    \item Kans op underfitting bij een ongeschikte parameterinstelling, wat
    resulteert in een hoge foutmarge op zowel trainings- als testdata.
\end{itemize}

\medskip

\paragraph{Autoencoder}


Een autoencoder is een neuraal netwerk bestaande uit 2 componenten: een encoder en een decoder. Deze twee componenten zijn verbonden door een tussenlaag genaamd latent space. 
De encoder comprimeert de input in het model naar een compacte representatie in deze latent space, dit zorgt ervoor dat enkel de meest essentiële kenmerken van de data worden behouden. 
Vervolgens probeert de decoder vanuit deze compacte representatie de oorspronkelijke output zo nauwkeurig mogelijk te reconstrueren. Het model wordt getrained de reconstructiefout, het verschil tussen originele input en gereconstrueerde output, te minimaliseren\autocite{Chikezie2025}.

Autoencoders worden in de context van anomaliedetectie voornamelijk op een semi-supervised manier geïmplementeerd. Doordat het model getraind wordt op normale input, is het instaat om gelijkaardige input te reconstrueren.
Echter wanneer input het model passeert dat eerder afwijkend is van het geleerde patroon, zal de reconstructiefout aanzienlijk hoger liggen doordat in de latent space niet de juiste representatie gevormd kon worden. Deze reconstructiefout is de basis als indicator voor een mogelijk anomalie\autocite{Chikezie2025}. 


Een datapunt wordt als anomalie beschouwd wanneer de reconstructiefout de vooraf bepaalde drempelwaarde overschrijdt. Een te lage drempelwaarde kan leiden tot een overvloed aan false positives, terwijl een te hoge waarde mogelijks anomalieën zal missen\autocite{Angiulli2023}.


\textbf{Voordelen:}
\begin{itemize}
    \item Breed toepasbaar en populair in diverse domeinen zoals netwerkdetectie,
    industriële inspectie en beeldanalyse, wat aantoont dat autoencoders flexibel
    inzetbaar zijn voor anomaliedetectie.
    \item In staat om complexe, niet-lineaire patronen in data te modelleren,
    waardoor ze geschikt zijn voor hoog-dimensionale en complexe datasets.
\end{itemize}

\textbf{Nadelen:}
\begin{itemize}
    \item Mogelijkheid tot \textit{out-of-bounds reconstruction}, waarbij
    datapunten ver buiten de trainingsdata toch een lage reconstructiefout
    krijgen, wat leidt tot false negatives.
    \item Gevoeligheid voor generalisatie-effecten zoals interpolatie en
    extrapolatie, waardoor anomalieën kunnen overlappen met normale data en
    onopgemerkt blijven\autocite{Bouman2025}.
    \item De keuze van de netwerk-architectuur, waaronder het aantal lagen en de dimensie van de bottleneck-laag, vereist experimentatie en afstemming. Een te grote bottleneck kan ertoe leiden dat het model ook afwijkende data goed reconstrueert, waardoor anomalieën worden gemist \autocite{Chikezie2025}.
\end{itemize}
\clearpage
\begin{table}[H]
    \centering
    \caption{Vergelijking van machine learning modellen voor anomaliedetectie}
    \label{tab:ml-modellen-vergelijking}
    \begin{tabular}{p{3cm} p{2.8cm} p{4.5cm} p{4.5cm}}
    \toprule
    \textbf{Model} & \textbf{Categorie} & \textbf{Voordelen} & \textbf{Nadelen} \\
    \midrule
    Random Forest 
      & Supervised 
      & Robuust tegen overfitting; geen expliciete featureselectie nodig 
      & Lage interpreteerbaarheid; prestaties dalen bij sterk gecorreleerde features \\
    \addlinespace
    Support Vector Machine 
      & Supervised 
      & Sterk in hoog-dimensionale data; capabel voor klassengrenzen 
      & Lage interpreteerbaarheid; trainingstijd stijgt exponentieel bij grote datasets \\
    \addlinespace
    Isolation Forest 
      & Unsupervised 
      & Lineaire tijdcomplexiteit; effectief door directe isolatie van anomalieën 
      & Moeilijkheden bij overlappende clusters; gevoelig voor dimensionaliteit \\
    \addlinespace
    Local Outlier Factor 
      & Unsupervised 
      & Detecteert lokale anomalieën; geen gelabelde data vereist 
      & Gevoelig voor parameterkeuze (k); kans op false positives aan clusterranden \\
    \addlinespace
    One-Class SVM 
      & Semi-supervised 
      & Geschikt voor één klasse; goede prestaties bij imbalanced datasets 
      & Sterke afhankelijkheid van hyperparameters; kans op underfitting \\
    \addlinespace
    Autoencoder 
      & Semi-supervised 
      & Breed toepasbaar; modelleert complexe niet-lineaire patronen 
      & Risico op out-of-bounds reconstructie; gevoelig voor generalisatie-effecten \\
    \bottomrule
    \end{tabular}
\end{table}
\clearpage

\subsection{Het belang van tijd in anomaliedetectie}%
\label{subsec:belang-tijd}
Tijd speelt een centrale rol in anomaliedetectie voor web- en api verkeer. Aanvallen zoals DDoS, credential stuffing en API-misbruik zullen zelden te detecteren zijn op basis van één enkele request. Deze aanvallen zullen eerder pas zichtbaar zijn voor een model wanneer ze worden geanalyseerd op een bepaalde tijdsspanne\autocite{Schummer2024}.%  Machine Learning-Based Network Anomaly Detection https://www.researchgate.net/publication/387151644_Machine_Learning-Based_Network_Anomaly_Detection_Design_Implementation_and_Evaluation
Om gedragsmatige patronen te kunnen detecteren wordt ruwe logdata samengevoegd over tijdsvensters. Dit is wanneer requests van eenzelfde client, endpoint of sessie samengebracht worden binnen een interval, dit kan bijvoorbeeld tussen de 1 en 5 minuten zijn. Per tijdsvenster wordt vervolgens de relevante features geanalysseerd\autocite{Sowmya2023}. % https://www.researchgate.net/publication/372514849_API_Traffic_Anomaly_Detection_in_Microservice_Architecture

De keuze van tijdsvenster is afhankelijk van de data waarop het ML model gebruikt zal worden, doorgaans zijn er twee mogelijkheden: tumbling windows en sliding windows. Bij tumbling windows worden er gebruik gemaakt van niet overlappende intervallen. Bij sliding windows is er sprake van een continue verschuiving van het interval om zo een gelijkmatiger beeld te geven van de gedragsverandering. Voor real-time detectie is sliding window aanbevolen doordat deze sneller reageren op plotse gedragswijzigingen\autocite{Schummer2024}. % Machine Learning-Based Network Anomaly Detection https://www.researchgate.net/publication/387151644_Machine_Learning-Based_Network_Anomaly_Detection_Design_Implementation_and_Evaluation

\subsection{Concept drift in anomaliedetectie}%
\label{subsec:concept-drift}

Machine-learningmodellen in de context van anomaliedetectie worden getrained op data die het normale verkeersgedrag van op dat moment representeert. In werkelijkheid verandert dit normaal gedrag continu: het lanceren van nieuwe features, het stijgen of dalen van verkeer van gebruikers en aanvallers die zich aanpassen. Deze verschuiving van eigenschappen en gedrag in data wordt ook wel concept drift genoemd. Wanneer concept drift optreed zonder dat het model wordt aangepast, ontstaat er een groeiende kloof tussen wat het model als normaal beschouwt en wat in de werkelijkheid op dat moment normaal is, het gevolg hiervan is een daling in detectieperformantie.   %https://dl.acm.org/doi/epdf/10.1145/3152494.3152501


\section{Feature selectie en engineering voor anomaliedetectie}
\label{sec:ML-methoden-anomalie}
\subsection{Belang van data quality}
\label{subsec:toegevoegde-waarde-ML}
% Uitleggen waarom goede data data essentieel is voor ML success 
Anomaliedetectie binnen Web- en API-beveiliging is sterk afhankelijk van de beschikbare data. Slechte data kwaliteit zoals onvolledige, inconsistente API of web metrics zorgen voor onnauwkeurige resultaten en verhogen het risico op false positives en negatives. Het verzekeren van data integriteit door correcte collectie, opslag en preprocessing is cruciaal voor effectieve anomalie detectie \Autocite{Meegle2025}.

\subsection{Van ruwe logdata naar gedragsprofielen}
\label{subsec:log-gedrags}
% Lijst de belangrijkste features voor anomaly detection in Web en API-context
Ruwe logvelden zoals IP-adres, statuscode of referrer vertellen op zichzelf weinig of een client zich al dan wel of niet normaal gedraagt. Een enkele request dat een 404 statuscode heeft is op zich niet verdacht, maar tientallen 404 requests van hetzelfde ip binnen een tijdsframe van enkele minuten kan wijzen op scanning. 
De meerwaarde voor een machine learning model ontstaat pas wanneer individuele logvelden worden samengevoegd en afgeleide kenmerken worden berekend die een gedragspatroon kunnen beschrijven over een tijdsvenster.
Dit proces wordt beschreven als feature engineering en is een belangrijke stap in anomalie detectie, zonder betekenisvolle features is het voor een goed presterend algoritme onmogelijk om een onderscheid te kunnen maken tussen normaal of afwijkend verkeer.%chandola en sowmya
Aangezien web- en API-verkeer verschillende kenmerken vertonen en andere aanvalspatronen kennen, spiltsen we binnen dit onderzoek de relevante features op in twee groepen. 
\subsubsection{API-specifieke features}
Voor het anomaliedetectie model voor API-verkeer, wordt er vertrokken van een set ruwe logvelden die rechtstreekse relevant zijn aan de API-endpoints: client IP-adres, API-endpoint, response data size, HTTP-statuscode, timestamp en request type. Deze kenmerken worden niet rechtstreeks gebruikt als input voor het model maar, dienen eerst te worden geaggregeerd over een tijdsvenster en dit expliciet per client-IP~\autocite{Sowmya2023}. 

De afgeleide features worden onderverdeeld in drie categorieën: volumetrische kenmerken, responsekenmerken en endpoint gedragskenmerken. 

\textbf{Volumetrische kenmerken} beschrijven de intensiteit waarop een client de API endpoints aanspreekt. Het totale aantal API-calls per IP binnen het tijdsvenster is de meest directe aanwijzing dat een IP mogelijks bezig is met flooding van het endpoint of geautomatiseerd misbruik \autocite{Sowmya2023}. Wanneer een IP binnen een kort tijdsvenster een afwijkend groot datavolume ontvangt, kan dit ook wijzen op het systematisch ophalen van gevoelige gegevens via de API \autocite{Sowmya2023, OWASP2023}.

\textbf{Responsekenmerken} deze kenmerken beschrijven het resultaat van uitgevoerde requests richting de API. Een verhoogd aantal 4xx-statuscodes kunnen een aanwijzing zijn van verkennend gedrag waarbij een IP systematisch endpoints bevraagt\autocite{OWASP2023}.

\textbf{Endpoint gedragskenmerken} beschrijven hoe een IP de endpoints benadert. Hierbij zijn de kernfeatures het aantal aangesproken API-endpoints en het totale aantal requests. Een afwijkend volume kan wijzen op endpoint abuse of brute force-aanvallen\autocite{Sowmya2023}. Een hoge verhouding met veel unieke endpoints in een tijdsvenster kan wijzen op het systematisch in kaart brengen van het API-oppervlakte \autocite{OWASP2023}.

\subsubsection{Web-specifieke features}
Voor het model dat anomalieën detecteert in HTTP web-verkeer, wordt vertrokken vanaf basis webserverlogvelden zoals: source IP, timestamp, method, URI, referrer. De gegevens worden gegroepeerd per source IP binnen een tijdsvenster, hierna worden voor dit specifieke ip relevante afgeleide features berekend die het gedrag van de client kunnen representeren~\autocite{Chua2024}. 

De afgeleide features worden onderverdeeld in vier categorieën: frequentiekenmerken, inhoudskenmerken, timing-kenmerken en gedragskenmerken. Deze categorisatie sluit aan bij de aanbevelingen uit de literatuur om zowel intensiteit als aard van het clientgedrag te onderzoeken \autocite{Chua2024, Schummer2024}.

\textbf{Frequentiekenmerken} beschrijven hoe intensief en op welke manier een client het systeem gebruikt binnen een bepaald tijdsvenster. Deze kenmerken worden berekend per IP-adres en geven inzicht in het algemene gebruikspatroon van een client.
De belangrijkste kenmerken zijn het totale aantal requests per IP en de verdeling van HTTP-statuscodes. Een afwijkend hoge frequentie aan requests kan wijzen op scraping of een DDoS-achtig patroon, terwijl hoge 4xx-responses kan duiden op scanning of brute force-pogingen\autocite{Chua2024, OWASP2023}.

\textbf{Inhoudskenmerken} geven een beeld over de requests zelf. Afwijkende requests bevatten vaak een afwijkend patroon zoals lange URI's die kunnen wijzen op injection-pogingen. Relevante kenmerken kunnen de gemiddelde en maximale URI-lengte zijn en de frequentie van specifieke URI-patronen binnen één IP\autocite{Chua2024}. 

\textbf{Timingkenmerken} beschrijven het tijds patroon van de requests. Normaal verkeer zal doorgaans variatie vertonen in de intervallen tussen requests, terwijl scripts en geautomatiseerde tools eerder uniforme intervallen vertonen. Daarom zijn de gemiddelde tijd en de standaardafwijking indicatieve features voor het onderscheid tussen menselijk of geautomatiseerd \autocite{Schummer2024, Benova2024}. 

\textbf{Gedragskenmerken} representeren de consistentie en diversiteit in het gedrag van de client. Voorbeelden hiervan zijn het gebruik van verschillende HTTP-methoden, verhouding tussen unieke paden, totale aantal requests en consistentie in de user-agent. Normaal verkeer zal zich doorgaans beperken tot Get en POST methoden en zal zelden in een tijdsvenster wisselen van user-agent \autocite{Benova2024, Chua2024}.  

\section{Evaluatie en performance metrics}
\label{sec:ML-methoden-anomalie}
\subsection{Evaluatie bij supervised anomaliedetectie}
\label{subsec:toegevoegde-waarde-ML}
Bij supervised anomaliedetectie beschikt de onderzoeker over ten minste één volledige gelabelde dataset waarbij elk punt expliciet geclassificeerd is als normaal of abnormaal. Dit maakt het mogelijk om de voorspellingen van het model rechtstreeks te vergelijken met de bekende grondwaarde en klassieke classificatiemetrics te berekenen \autocite{Goldstein2016}.

Precision meet het aandeel correct gedetecteerde anomalieën ten opzichte van alle als anomalie geclassificeerde observaties. Recall meet het aantal correcte anomalieën ten opzichte van het totale aantal dat zich in de dataset bevindt. Beide metrics zijn belangrijk in de evaluatie van anomaliedetectie aangezien beide zich focussen op de minderheidsklasse binnen het model (de anomalieën) en zo een nauwkeuriger beeld bieden over de detectiekwaliteit. De F1-score combineert deze twee metrics als gemiddelde, waardoor de F1-score enkel kan stijgen indien beide waarden stijgen\autocite{Chadha2020}.  

\subsection{Evaluatie bij unsupervised anomaliedetectie}
\label{subsec:toegevoegde-waarde-ML}
% Definiëren van de metrics die je gaat gebruiken voor de evaluatie van het model. 
Bij unsuperised anomaliedetectie is een betrouwbare gelabelde dataset of grondwaarde doorgaans niet beschikbaar, dit is vaak ook de reden waarom in eerste instantie gekozen is voor de implementatie van een unsupervised model. Dit heeft echter wel rechtstreeks gevolgen op de evaluatie van het model, metrics zoals recall of F1-score vereisen kennis van het totale aantal werkelijke anomalieën die zich in de dataset bevinden \autocite{Goldstein2016}.

Om toch de detectiekwaliteit van deze modellen te kunnen beoordelen worden voor deze modellen vaak alternatieve strategieën gebruikt. Een veelgebruikte strategie is handmatige validatie, waarbij de detecties gevonden door een model worden beoordeeld door een domeinexpert of door de onderzoeker zelf \textcite{Chua2024} pasten deze aanpak toe door de testdata van hun Isolation Forest-model handmatig te labelen, waarna precision, recall en F1-score berekend konden worden.

\textcite{Rochner2023} gebruikten een vergelijkbare strategie bij unsupervised anomaliedetectie in medische dossiers, waarbij domeinexperts een steekproef van de detecties beoordeelden om de precision van het model te kunnen bepalen. 

In de praktijk worden meerdere technieken vaak gecombineerd om vanuit meerdere invalshoeken een beeld te creeëren over hoe het model detecteert, de keuze hiervan is vaak afhankelijk van de beschikbare methoden. 

\section{Conclusie literatuurstudie}
\label{sec:ML-methoden-anomalie}

De literatuurstudie toont aan dat traditionele detectiesystemen zoals WAF's en signature based IDS-systemen effectief zijn in het herkennen van eerder geziene en gedocumenteerde dreigingen, maar beperkingen kunnen vertonen bij het detecteren van ongeziene of subtiele aanvalsvormen. Toenemende automatisering en complexiteit van aanvallen versterkt deze problematiek. 

Uit de vergelijking van machine learning-benaderingen blijkt dat unsupervised methoden het meest geschikt zijn wanneer gelabelde trainingsdata ontbreekt, wat in de meeste operationele security-omgevingen het geval is. 

Isolation Forest onderscheidt zich binnen deze categorie door zijn lineaire tijdcomplexiteit, lage geheugenverbruik en de eigenschap dat het anomalieën rechtstreeks isoleert in plaats van eerst normaal gedrag te modelleren. Dit maakt het bijzonder geschikt voor toepassing binnen een real-time detectiesysteem dat grote hoeveelheden logdata moet verwerken.

Daarnaast blijkt feature engineering een bepalende invloed te hebben op de detectiekwaliteit van het model. Ruwe logvelden zeggen weinig op zich, de meerwaarde ontstaat door deze velden te aggregeren over een tijdsvenster en per client-IP, hierdoor worden gedragspatronen van gebruikers zichtbaar en kunnen afwijkingen worden gedetecteerd. 

Deze inzichten vormen de basis voor de modelkeuze, feature engineering en architectuur van het proof-of-concept dat in de volgende hoofdstukken wordt uitgewerkt.
% Voorbeeldtekst van het sjabloon  



